\documentclass{memoir}
\usepackage{amssymb}
\usepackage{tcolorbox}
\usepackage{amsmath}
\usepackage[italian]{babel}
\usepackage[T1]{fontenc}        % Codifica dei font migliorata
\usepackage[utf8]{inputenc}     % Supporto caratteri UTF-8 (di default in LaTeX moderni)
\usepackage{ulem}
\usepackage{tikz}
% \usepackage{unicode-math}
%\setmathfont{Garamond Math}
%\\usepackage[varbb]{newpxmath}
%\usepackage[activate={true,nocompatibility},final,tracking=true,kerning=true]{microtype}
% microtype makes text look nicer
\usepackage[italian]{babel}     % Lingua italiana
\title{Algebra Lineare}
\usepackage{tcolorbox}
\usepackage{dirtytalk}

\tcbset {
  base/.style={
    arc=0mm,
    bottomtitle=0.5mm,
    boxrule=0mm,
    colbacktitle=black!10!white,
    coltitle=black,
    fonttitle=\bfseries,
    left=2.5mm,
    leftrule=1mm,
    right=3.5mm,
    title={#1},
    toptitle=0.75mm,
  }
}


\definecolor{brandblue}{rgb}{0.34, 0.7, 1}
\newtcolorbox{mainbox}[1]{
  colframe=brandblue,
  base={#1}
}

\newtcolorbox{subbox}[1]{
  colframe=black!30!white,
  base={#1}
}
\newtheorem{definizione}{\textbf{Definizione}}[section]
\newtheorem{esempio}{Esempio}[section]
\begin{document}
\maketitle
%% qualche problemino di organizzazione con i capitoli: metto vettori e matrici qui? prima di sistemi lineari

\chapter{Punti, Vettori e matrici}
In algebra lineari gli attori principali che studiamo sono vettori e matrici. Seguono le definizioni strettamente necessarie per comprendere i prossimi capitoli.
\section{Punti e vettori}
Un elemento di \(\mathbb{R}^n\) è una lista ordinata di numeri, ma questa lista può essere interpretata in 2 modi: come un punto o come un vettore.
\begin{tcolorbox}
  \begin{definizione}
    punti e vettori
  \end{definizione}
  L'elemento di \(\mathbb{R}^n\) con coordinate \(x_1,x_2,\cdots, x_n\) può essere interpretato come il \textit{punto} \(\textbf{x}= \begin{pmatrix}
    x_1\\
    x_2\\
    \vdots\\
    x_n
 \end{pmatrix}
 \) o il vettore \(\vec{x}=\begin{bmatrix}
   x_1\\
   x_2\\
   \vdots\\
   x_n
 \end{bmatrix}
 \), che rappresentano un incremento rispetto un punto in \(\mathbb{R}^n\)
\end{tcolorbox}

\begin{esempio}Differenza tra punti e vettori \end{esempio} L'elemento di \(\mathbb{R}^n\) con coordinate x=2 e y=3 può essere interpretato come il punto \(\begin{pmatrix}
    2\\
    3

  \end{pmatrix}\) o come il vettore \(\begin{bmatrix}
    2\\
    3
  \end{bmatrix}
  \) che significa:\say{Inizia da qualsiasi punto e spostati di due unità a destra e tre unità in alto}
  \subsection{I vettori possono essere sommati, i punti no}
  Possiamo identificare la posizione di una città su una mappa, cioè le sue coordinate, ma non aggiungiamo mai le loro posizioni. Ad esempio non aggiungiamo mai le posizioni di \textbf{Roma} o \textbf{Milano}. Ha senso invece la loro differenza (distanza): quanto dobbiamo aggiungere a un punto per arrivare all'altro. La differenza tra un punto \textbf{a} e un punto \textbf{b} è il vettore  \(\vec{\textbf{a}-\textbf{b}}\).

  \begin{definizione}
    Somma tra vettori
  \end{definizione}
  I vettori sono sommati aggiungendo le loro componenti:
  \begin{equation}
    \begin{bmatrix}
      v_1\\
      v_2\\
      \vdots\\
      v_n
    \end{bmatrix} +
    \begin{bmatrix}
      w_1\\
      w_2\\
      \vdots\\
      w_n

    \end{bmatrix} = \begin{bmatrix}
      v_1+w_1\\
      v_2+w_2\\
      \vdots\\
      v_n+w_n
    \end{bmatrix}
    \end{equation}

    Nel piano, la somma \(\vec{v}+\vec{w}\) è la diagonale del parallelogramma i cui 2 lati adiacenti sono \(\vec{v}\) e \(\vec{w}\)

    \section{prodotto tra scalare e vettore}
    Il prodotto tra scalare e vettore è molto semplice: ogni elemento del vettore è moltiplicato per lo scalare.

    \begin{equation}
      a \begin{bmatrix}
        x_1\\
        \vdots\\
        x_n
      \end{bmatrix} = \begin{bmatrix}
        ax_1\\
        \vdots\\
        ax_n
      \end{bmatrix}
     3 \cdot \begin{bmatrix}
    2\\
    4\\
    \sqrt(2)
  \end{bmatrix} = \begin{bmatrix}
    6\\
    12\\
    3\sqrt{2}
  \end{bmatrix}
\end{equation}
\section{Matrici}
\begin{tcolorbox}
  Una matrice  \( m \times n \) è una lista rettangolare di altezza \(m\) e \(n\) lunga. Denotiamo \textit{l}' insieme delle matrici con Mat(m,n) o con \(M_{m,n}(\mathbb{K})\) in campo \(\mathbb{K}\). Il numero \(a_{ij}\) è l'elemento di posto (i,j) della matrice A, dove il primo è l'indice di riga, mentre il secondo è l'indice j è l'indice di colonna.
\end{tcolorbox}
Usiamo lettere maiuscole per denotare le matrici. Un vettore \(\vec{v} \in \mathbb{R}^m\) è una matrice \(m \times 1\); un numero è una matrice \(1 \times 1\). Ma perchè dobbiamo occuparci di matrici invece di matrici quando già esistono i vettori? la risposta che le matrice ci permettono di fare una operazione: la moltiplicazione tra matrici.
\subsection{moltiplicazioni tra matrice}
Il primo elemento della riga AB è ottenuta moltiplicando gli elementi della riga per gli elementi della corrispondente colonna e poi aggiungendo quei prodotti assieme. Ad esempio il primo elemento della matrice AB è ottenuto così: \((2 \times 1) + (-1 \times 3)= -1\)
\begin{equation}
  \begin{bmatrix}
    2 & -1\\
    3 & 2
  \end{bmatrix} \begin{bmatrix}
    1 & 4 & -2\\
    3 & 0 & 2
  \end{bmatrix} =
  \begin{bmatrix}
    -1 & 8 & -6\\
    9 & 12 & -2
    \end{bmatrix}
  \end{equation}

  \begin{tcolorbox}
    Se A è una matrice m × n
il cui (i, j)-esimo elemento è \(a_{ij}\), e B è una matrice n × p il cui (i, j)-esimo elemento
è bi,j , allora C = AB è la matrice m × p con elementi:
\begin{equation}
c_{i,j} =   \sum_{k=1}^n {a_{i,k}b_{k,j}} = a_{i,1}b_{1,j}+a_{i,2}b_{j,2} + \cdots +a_{i,n}b_{n,j}
  \end{equation}
\end{tcolorbox}
\subsection{Matrci speciali}
\subsection{Matrice identità}
denotata con \(I_n\) è una matrice quadrata di dimensione \(n\) dove tutti gli elementi della diagonale principale sono costiuti dal numero 1, mentre i restanti elementi sono 0.
\begin{equation}
I_n = \begin{bmatrix}
  1 & 0 & \cdots & 0\\
  0 & 1 & \cdots & 0\\
  \vdots & \vdots & \ddots\\
  0 & 0 & \cdots & 1
\end{bmatrix}
\end{equation}

\chapter{Sistemi Lineari}
\section{sistemi di equazioni}
Diciamo che vogliamo risolvere questo sistema di equazioni:

\begin{align*}
x_1+x_2 + x_4 + x_5&= 2\\
2x_1+2x_2-x_3-x_5&=-1 \\
3x_1+3x_2-2x_3-x_4-x_5 &= -2
\end{align*}

  Possono essere utilizzate diverse tecniche, ma la più usata è \emph{l'algoritmo di Gauss}. Prima di tutto si scrive il sistema in forma matriciale:
 \begin{align}
  \underbrace{
  \begin{bmatrix}
    1 & 0 &0 & 1 &1 \\
    2 & 2 &-1 & 0 &-1\\
    3 & 3 & -2 & -1 &-1
  \end{bmatrix}
    }_{\text{matrice coefficenti A}}
    \quad
  \underbrace{
  \begin{bmatrix}
    x \\
    y \\
    z
  \end{bmatrix}}_{\text{vettore incognite } \vec{x}}
    \quad
    \underbrace{
    \begin{bmatrix}
      2 \\
      -1 \\
      -2
    \end{bmatrix}
    }_{ \text{vettore termini noti } \vec{b} }
    \end{align}
    Più generalmente il sistema di equazioni \(m\) in \(n\) incognite è scritto come
    \[
      \begin{matrix}
        a_\text{1,1}x_1 & + \dots + &x_\text{1,n}x_n &= b_1 \\
        \vdots& \cdots & \vdots & \vdots\\
        a_\text{m,1}x_1 & +\dots+&a_\text{m,n}x_n &= b_m
        \end{matrix}
      \]
      ed uguale ad \(A\vec{x}=\vec{b}\):
      \begin{align*}
      \begin{bmatrix}
        a_\text{1,1} & \dots & a_\text{1,n}\\
        \vdots  & & \vdots \\
        a_\text{m,1} & \dots & a_\text{m,n}
      \end{bmatrix}
      \begin{bmatrix}
          x_1\\
          \vdots\\
          x_n
      \end{bmatrix}
        =
        \begin{bmatrix}
          b_1\\
          \vdots\\
          b_m
          \end{bmatrix}
      \end{align*}
      Con \( [A|\vec{b}]\), o anche \(A'\), si rappresenta \textbf{la matrice completa del sistema}. La matrice completa contiene sia coefficenti e termini noti. Tutte le operazioni per risolvere un sistema sono fatte sulla matrice completa del sistema. Una soluzione del sistema è \textit{n-upla} \(v_1,\dots, v_n\) che se sostituita all'interno del vettore \(\vec{x}\) soddisfano le equazioni del sistema. Il sistema è compatibile se esiste almeno una soluzione. Un soluzione del sistemi è una n-pla \(\begin{bmatrix}
      x_1\\
      x_2\\
      \vdots\\
        x_m\
        \end{bmatrix}\) che se sostituita rende soddisfate tutte le equazioni

      \section{Operazioni che lasciano il sistema invariato}
      Si può risolvere un sistema di equazioni lineari \(A\vec{x}=\vec{b}\) utilizzando \textbf{L'algoritmo di Gauss} e usando le operazioni che lasciano il sistema invaritato.

      \begin{tcolorbox}[
        title=Matrice,
  colback=blue!5!white,
  colframe=blue!75!black,
  fonttitle=\bfseries
  {
    Una matrice è una tabella rettangolare di m righe con n colonne. Una matrice quadra ha lo stesso numero di righe e colonne.

  }
\end{tcolorbox}

  \begin{tcolorbox}[
  title=Operezioni di riga,
  colback=blue!5!white,
  colframe=blue!75!black,
  fonttitle=\bfseries
  ]
  \enumerate
  \item moltiplicare una riga per un numero diverso da 0
  \item  aggiungere una riga ad un altra riga
  \item scambiare due righe

\end{tcolorbox}
Le operazioni di riga sono importanti perchè sono alla base dell' algoritmo di gauss. Anche perchè dato che si basano su operazioni aritmetiche elementari possono essere eseguite facilmente da un computer.

\subsection{L'algoritmo di Gauass}
I sistemi lineari possono essere portati in una forma equivalente che ci permette di capire facilmente se sono risolvibili: \textbf{La forma a scala}.
\begin{tcolorbox}[title=sistemi equivalenti]
  Due sistemi si dicono equivalenti se hanno le stesse soluzioni.
\end{tcolorbox}
\(\mathit{}\)
\subsection{La forma a scala}
La ragione per cui ci interessano tanto i sistemi a scala è che rendo motlo facile risolvere i sistemi lineari.
\begin{definizione}
  Matrice a scala
\end{definizione}
  Una matrice si dice a scala se soddisfa le seguenti proprietà:
  \begin{enumerate}
    \item Se una riga è nulla, allora tutte le righe sotto sono nulle.
          \item Sotto il primo elemento non nullo di ciascuna riga, e sotto tutti gli zeri che lo precedono, ci
sono elementi nulli. Questo elemento si chiama pivot.
  \end{enumerate}
\begin{equation}
\begin{pmatrix}
\uline{p_1} & * & *& *\\
0 & \uline{p_1} & * & * \\
0 & 0 & \uline{p_1} & * \\
\vdots & \vdots & \vdots & \vdots\\
  0 & 0 & 0 & 0
\end{pmatrix}
\end{equation}
\subsection{matrice a scala ridotta}
La matrice a scala ridotta differisce dalla matrice a scala solo per i pivot uguali ad 1.
\begin{equation}
\begin{pmatrix}
\uline{1} & * & *& *\\
0 & \uline{1} & * & * \\
0 & 0 & \uline{1} & * \\
  \vdots & \vdots & \vdots & \vdots\\
  0 & 0 & 0 & 0\\
\end{pmatrix}
\end{equation}
Se \~A contiene un gradino nei termini noti allora non ha una soluzione. Se non c'è un gradino nella colonna dei termini noti allora il sistema ha la forma:
\begin{equation*}
  \begin{bmatrix}
    3 & 2 & 0 & 5\\
    0 & 0 &-1 & -3\\
    0 & 0 & 0 & 0
  \end{bmatrix} \Longrightarrow \left\{\begin{aligned}
    3x+2y&=5\\
    -z&=-3
  \end{aligned}\right. \Longrightarrow
\left\{\begin{aligned} x&=\frac{1}{3}(5-2y)\\
  z&=3
  \end{aligned}\right.
\end{equation*}
\subsection{Rango di una matrice e numero di soluzione}
L' algoritmo di Gauss ci permette di portare una matrice completa del sistema in una forma equivalente in cui è facile risolvere le equazioni. Una volta portata in forma a scalini possiamo partire dall'ultima equazione e sostituire all'indietro.
\subsection{compatibilità}
Un sistema è compatibile se esiste almeno una soluzione. Questo accade quando il numero dei pivot è uguale al numero delle incognite, in altre parole la lunghezza dei gradini è 1.
\subsection{non compatibilità}
Se l'ultimo pivot cade nella colonna dei termini noti allora il sistema non soluzione.
\subsection{soluzioni inifinite con parametri liberi}
Quando nella matrice in alcune colonne non ci sono pivot allora questo vuol dire che la colonna che corrisponde a quella variabile è un parametro libero chiamato \(t\) che può assumere un valore a nostro piacimento. Questo avviene se il numero di pivot è inferiore al numero di incognite.
\begin{tcolorbox}[title=rango]
  Il rango di una matrice è il numero di pivot nella matrice a scalini complete.
  \begin{equation}
    rg(A) =  \text{\# numero di pivot}
    \end{equation}
  \end{tcolorbox}
\end{document}
